\preface % Note: \preface JANGAN DIHAPUS!

Segala puji dan syukur kehadiran Allah SWT yang telah melimpahkan rahmat
dan hidayah-Nya kepada kita semua, sehingga penulis dapat menyelesaikan penulisan tugas akhir yang berjudul “Penerapan Algoritma K-Nearest Neighbors dalam Prediksi Kualitas Pengguna pada Aplikasi E-Learning Berbasis Web” yang telah dapat diselesaikan sesuai rencana. Penulis banyak mendapatkan berbagai pengarahan, bimbingan, dan bantuan dari berbagai pihak. Oleh karena itu, melalui tulisan ini penulis mengucapkan rasa terima kasih kepada:

\begin{enumerate}
    \item Orang tua tercinta Ayahanda XX dan Ibunda XX, ...
    
    \item Dr. Nizamuddin, S.Si., M.Info.Sc., selaku Ketua Departemen Informatika Fakultas MIPA Universitas Syiah Kuala.
    
    \item Bapak/Ibu XXxx,  selaku Dosen Pembimbing I, memberikan bimbingan dan arahan yang sangat dibutuhkan oleh penulis, yang akhirnya membantu penulis menyelesaikan tugas akhir ini.
    
    \item Bapak/Ibu XXXX, selaku Dosen Pembimbing II, memberikan bimbingan dan arahan yang sangat dibutuhkan oleh penulis, yang akhirnya membantu penulis menyelesaikan tugas akhir ini.
    
    \item Bapak/Ibu XXXX selaku dosen pembahas yang telah memberikan masukan dan saran yang sangat berharga bagi penulis dalam menyelesaikan tugas akhir ini.
    
    \item  Bapak/Ibu XXXX, selaku dosen pembimbing dan dosen wali yang telah memberikan bimbingan, masukan, dan saran yang sangat berharga bagi penulis dalam menyelesaikan tugas akhir ini. 
    
    \item Seluruh Dosen dan Staf di Jurusan Informatika Fakultas MIPA, atas ilmu dan bimbingan yang diberikan selama perkuliahan, 
    
    \item Tambahan lainnya...
    
    \item Sahabat dan teman-teman seperjuangan dari Jurusan Informatika USK 2021 lainnya, telah memberikan dukungan selama penulis menempuh pendidikan di jurusan Informatika USK.
\end{enumerate}

%\vspace{1.5cm}

Penulis mengakui adanya kekurangan dalam tulisan ini, baik dari aspek materi, metode, maupun bahasa yang digunakan. Oleh karena itu, penulis mengundang kritik dan saran yang konstruktif dari para pembaca untuk meningkatkan kualitas Tugas Akhir ini. Penulis berharap bahwa tulisan ini dapat memberikan manfaat bagi banyak orang dan berkontribusi pada kemajuan ilmu pengetahuan.

\vspace{1cm}

\begin{tabular}{p{7.5cm}l}
	&Banda Aceh, 03 Maret 2025\\
	&\\
	&\\
	&\multirow{1.5}{7.5cm}{\underline{Nama Nama}} \\ 
	&NPM. 2X081070100XX \\
\end{tabular}
